\section*{Problem 1}

已知矩阵
\[
    A = \begin{bmatrix}
        1 & 1 \\
        1 & -1 \\
        1 & 0
    \end{bmatrix}.
\]
由于矩阵 $A$ 列满秩，其 Moore-Penrose 广义逆 $A^+$ 可以通过公式
\[
    A^+ = \left(A^\top A\right)^{-1}A^\top
\]
来计算。请写出详细的计算过程，求出 $A^+$。

\begin{proof}
    首先计算
    \[
        A^\top
        = \begin{bmatrix}
            1 & 1 & 1 \\
            1 & -1 & 0
        \end{bmatrix},
    \]
    从而
    \[
        A^\top A
        = \begin{bmatrix}
            1 & 1 & 1 \\
            1 & -1 & 0
        \end{bmatrix}
        \begin{bmatrix}
            1 & 1 \\
            1 & -1 \\
            1 & 0
        \end{bmatrix}
        = \begin{bmatrix}
            3 & 0 \\
            0 & 2
        \end{bmatrix}.
    \]
    因此
    \[
        \left(A^\top A\right)^{-1}
        = \begin{bmatrix}
            \frac{1}{3} & 0 \\
            0 & \frac{1}{2}
        \end{bmatrix}.
    \]
    代入列满秩矩阵的广义逆公式，得到
    \[
        \begin{aligned}
            A^+
            &= \left(A^\top A\right)^{-1}A^\top \\
            &= \begin{bmatrix}
                \frac{1}{3} & 0 \\
                0 & \frac{1}{2}
            \end{bmatrix}
            \begin{bmatrix}
                1 & 1 & 1 \\
                1 & -1 & 0
            \end{bmatrix} \\
            &= \boxed{\begin{bmatrix}
                \frac{1}{3} & \frac{1}{3} & \frac{1}{3} \\
                \frac{1}{2} & -\frac{1}{2} & 0
            \end{bmatrix}}.
        \end{aligned}
    \]
\end{proof}

\section*{Problem 2}

假设线性方程组 $A\bm{x} = \bm{b}$ 有解。利用广义逆的通解性质可知，该方程组的解可以统一表示为
\[
    \bm{x} = A^+\bm{b} + \left(I - A^+A\right)\bm{z},
\]
其中 $\bm{z}$ 为任意向量，且 $A^+$ 为 $A$ 的 Moore-Penrose 广义逆。请证明：
\[
    \bm{x}^* = A^+\bm{b}
\]
是该方程组所有解中二范数（模长）最小的解。

\begin{proof}
    因为方程组有解，可取 $\bm{x}_0$ 使得 $\bm{b} = A\bm{x}_0$。由 Moore-Penrose 广义逆的性质 $AA^+A = A$，有
    \[
        A\bm{x}^*
        = AA^+\bm{b}
        = AA^+A\bm{x}_0
        = A\bm{x}_0
        = \bm{b},
    \]
    因此 $\bm{x}^*$ 确实是方程组的一个解。

    记
    \[
        P = A^+A.
    \]
    由 Moore-Penrose 广义逆的性质可知
    \[
        P^\top = P, \qquad
        P^2 = A^+AA^+A = A^+A = P,
    \]
    因此 $P$ 是正交投影矩阵。又由 $A^+AA^+ = A^+$，有
    \[
        P\bm{x}^*
        = A^+AA^+\bm{b}
        = A^+\bm{b}
        = \bm{x}^*.
    \]
    对任意解
    \[
        \bm{x} = \bm{x}^* + (I - P)\bm{z},
    \]
    计算两个分量的内积可得
    \[
        \begin{aligned}
            \left(\bm{x}^*\right)^\top(I - P)\bm{z}
            &= \left(P\bm{x}^*\right)^\top(I - P)\bm{z} \\
            &= \left(\bm{x}^*\right)^\top P^\top(I - P)\bm{z} \\
            &= \left(\bm{x}^*\right)^\top\left(P - P^2\right)\bm{z}
             = 0.
        \end{aligned}
    \]
    因此 $\bm{x}^*$ 与 $(I - P)\bm{z}$ 正交，由勾股定理得到
    \[
        \|\bm{x}\|^2
        = \|\bm{x}^*\|^2 + \|(I - P)\bm{z}\|^2
        \ge \|\bm{x}^*\|^2.
    \]
    等号当且仅当 $(I - P)\bm{z} = \bm{0}$，即 $\bm{x} = \bm{x}^*$。故
    \[
        \boxed{\bm{x}^* = A^+\bm{b}}
    \]
    是该方程组所有解中唯一的最小二范数解。
\end{proof}

\section*{Problem 3}

已知 $A^+$ 是矩阵 $A$ 的 Moore-Penrose 广义逆，满足以下四条定义性质：
\[
    \begin{aligned}
        &\text{(1)}\quad AA^+A = A; &
        &\text{(2)}\quad A^+AA^+ = A^+; \\
        &\text{(3)}\quad \left(AA^+\right)^\top = AA^+; &
        &\text{(4)}\quad \left(A^+A\right)^\top = A^+A.
    \end{aligned}
\]
请利用上述四条性质，证明 $\left(A^+\right)^\top$ 也是 $A^\top$ 的 Moore-Penrose 广义逆，从而得出结论
\[
    \left(A^\top\right)^+ = \left(A^+\right)^\top.
\]

\begin{proof}
    令
    \[
        B = \left(A^+\right)^\top.
    \]
    对性质 (1) 和 (2) 分别取转置，得到
    \[
        A^\top B A^\top
        = \left(AA^+A\right)^\top
        = A^\top,
    \]
    以及
    \[
        B A^\top B
        = \left(A^+AA^+\right)^\top
        = B.
    \]
    再由性质 (4)，有
    \[
        A^\top B
        = A^\top\left(A^+\right)^\top
        = \left(A^+A\right)^\top
        = A^+A,
    \]
    所以
    \[
        \left(A^\top B\right)^\top
        = \left(A^+A\right)^\top
        = A^+A
        = A^\top B.
    \]
    同理由性质 (3)，有
    \[
        BA^\top
        = \left(A^+\right)^\top A^\top
        = \left(AA^+\right)^\top
        = AA^+,
    \]
    所以
    \[
        \left(BA^\top\right)^\top
        = \left(AA^+\right)^\top
        = AA^+
        = BA^\top.
    \]
    因此 $B$ 满足 $A^\top$ 的全部四条 Moore-Penrose 广义逆定义性质。由 Moore-Penrose 广义逆的唯一性，得到
    \[
        \boxed{\left(A^\top\right)^+ = B = \left(A^+\right)^\top}.
    \]
\end{proof}
